% --- IMPORTS ---
\documentclass[leqno]{article}
\usepackage{graphicx}
\usepackage{dsfont}
\usepackage{mathtools}
\usepackage{amssymb}
\usepackage{amsmath}
\usepackage{amsthm}
\usepackage{float}
\usepackage{ragged2e}
\usepackage{tensor}
\usepackage{fancyhdr}
\usepackage{empheq}
\usepackage[most]{tcolorbox}
\usepackage{enumitem}
\usepackage{dirtytalk}
\usepackage{marginnote}
\usepackage{microtype}
\usepackage[
  a4paper,
  left=0.8in,
  right=1in,
  top=1in,
  bottom=0.5in,
  headheight=14pt,
  headsep=0.4in
]{geometry}
\setlength{\parindent}{0cm}
% --- TikZ Packages ---
\usepackage{pgfplots}
\pgfplotsset{compat=1.18}
\usepgfplotslibrary{fillbetween, polar}
\usetikzlibrary{calc, positioning}
\usetikzlibrary{angles, quotes}
\usetikzlibrary{arrows.meta}
\usetikzlibrary{decorations.pathreplacing, decorations.markings}
\usetikzlibrary{patterns}
\usetikzlibrary{intersections}
% --- URL Packages ---
\usepackage[colorlinks=true,linkcolor=red,citecolor=magenta,urlcolor=black]{hyperref}%
\urlstyle{same}
% --- END OF IMPORTS ---

% --- CUSTOM COMMANDS ---
\RenewDocumentCommand{\marks}{o m}{%
  \IfValueTF{#1}%
    {\marginnote{\raggedright\textbf{[#2]}}[#1]}%
    {\marginnote{\raggedright\textbf{[#2]}}}%
}
\newcommand{\vspacerule}[2]{%
  \par\vspace*{#1}%
  \noindent\hrulefill\par
  \vspace*{#2}%
}
\definecolor{scarlet}{RGB}{205, 40, 75}
\newcommand{\questionitem}{%
  \item\phantomsection
  \addcontentsline{toc}{section}{Question \number\value{enumi}}%
}
\renewcommand{\contentsname}{Questions}
% --- END OF CUSTOM COMMANDS ---

% --- HEADER CONFIG ---
\pagestyle{fancy}
\fancyhead{}
\fancyfoot{}
\renewcommand{\headrulewidth}{0.9pt}
\lhead{Linear Transformations}
\chead{}
\rhead{\href{https://danieljsmith.org}{danieljsmith.org}}
% --- END OF HEADER CONFIG ---

\begin{document}
\vspace*{20pt}
\tableofcontents
\newpage
\begin{enumerate}[
  label=\textbf{\arabic*.},
  leftmargin=*,
  topsep=0pt,
  partopsep=0pt,
  parsep=0pt
]

% ---- Question 1 ----
\questionitem

\begin{enumerate}[label=\textbf{(\roman*)}]
  \item In each of the following cases, find a $2\times 2$ matrix that represents
  \begin{enumerate}[label=(\alph*)]
    \item a reflection in the line $y=x$ \marks{1}
    \item a rotation of $60^\circ$ clockwise about $(0,0)$ \marks{1}
    \item a rotation of $60^\circ$ clockwise about $(0,0)$ followed by a reflection in the line $y=x$ \marks{2} \\
  \end{enumerate}
  \item The triangle $T$ has vertices at the points $(1,2)$, $(5,5)$ and $(k,-1)$, where $k>-3$ is a constant. \\

  Triangle $T$ is transformed onto the triangle $T'$ by the matrix
  \[
  \begin{pmatrix}
  2 & 5 \\
  1 & -1
  \end{pmatrix}
  \]
  Given that the area of triangle $T'$ is $84$ square units, find the value of $k$. \marks{6} \\
\end{enumerate}

\newpage

% ---- Question 2 ----
\questionitem

\[
\mathbf{P}=\begin{pmatrix} p+1 & 2 \\ 2 & p+1 \end{pmatrix}
\] \\
where $p$ is a positive constant. \\

The matrix $\mathbf{P}$ represents a linear transformation $U$. \\
The parallelogram $S$ has vertices at the points with coordinates $(1,1)$, $(4,2)$, $(6,5)$ and $(3,4)$. \\
The area of the image of $S$ under the linear transformation $U$ is $35$. \\

\begin{enumerate}[label=\textbf{(\alph*)}]
  \item Determine the value of $p$. \marks{4} \\
\end{enumerate}

The transformation $V$ consists of a reflection in the line $y=x$ followed by a stretch of scale factor $2$ parallel to the $x$-axis with the $y$-axis invariant. The transformation $V$ is represented by the matrix $\mathbf{Q}$. \\

\begin{enumerate}[label=\textbf{(\alph*)}, resume]
  \item Write down the matrix $\mathbf{Q}$. \marks{2} \\
\end{enumerate}

Given that $U$ followed by $V$ is the transformation $W$, which is represented by the matrix $\mathbf{R}$ \\

\begin{enumerate}[label=\textbf{(\alph*)}, resume]
  \item find the matrix $\mathbf{R}$. \marks{2} \\
\end{enumerate}

\newpage

% ---- Question 3 ----
\questionitem

\begin{enumerate}[label=\textbf{(\alph*)}]
  \item Specify fully the transformation $\mathbf{T}$ of the plane associated with the matrix $\mathbf{M}$, where
  \[
  \mathbf{M} = \begin{pmatrix} \cos\theta & -\sin\theta \\ \sin\theta & \cos\theta \end{pmatrix}
  \]
  and $0<\theta<2\pi$. \marks{2} \\
  \item
  \begin{enumerate}[label=(\roman*)]
    \item Find $\det \mathbf{M}$. \marks{1}
    \item Deduce \textbf{two properties} of the transformation $\mathbf{T}$ from the value of $\det \mathbf{M}$. \marks{2} \\
  \end{enumerate}
  \item Prove that
  \[
  \mathbf{M}^n = \begin{pmatrix} \cos n\theta & -\sin n\theta \\ \sin n\theta & \cos n\theta \end{pmatrix}
  \]
  where $n$ is a positive integer. \marks{4} \\
  
  \item Hence specify fully a \textbf{single} transformation which is equivalent to $n$ applications of the transformation $\mathbf{T}$. \marks{1} \\
\end{enumerate}

\newpage

% ---- Question 4 ----
\questionitem

Consider the matrix
\[
A=\begin{pmatrix}
-\frac{1}{2} & \frac{\sqrt{3}}{2}\\[6pt]
\frac{\sqrt{3}}{2} & \frac{1}{2}
\end{pmatrix}
\] \\

\begin{enumerate}[label=\textbf{(\roman*)}]
  \item Calculate $A^2$. \marks{1} \\
\end{enumerate}

You are told that the matrix $A$ represents a reflection in a line through the origin. \\

\begin{enumerate}[label=\textbf{(\roman*)}, resume]
  \item Explain why your answer to part (i) is consistent with this information. \marks{1} \\
  \item By finding the invariant points of the transformation, determine the equation of the mirror line. \marks{3} \\
  \item Describe fully the transformation represented by
  \[
  B=\begin{pmatrix}
  -\frac{1}{2} & -\frac{\sqrt{3}}{2}\\[6pt]
  \frac{\sqrt{3}}{2} & -\frac{1}{2}
  \end{pmatrix}
  \]
  \marks[-30pt]{2}
  \item A composite transformation is formed by the transformation represented by $A$ followed by the transformation represented by $B$. Find the single matrix that represents this composite transformation. \marks{2} \\
  \item The composite transformation described in part (v) is equivalent to a single reflection. Find the equation of the mirror line of this reflection. \marks{1} \\
\end{enumerate}

\newpage

% ---- Question 5 ----
\questionitem

The transformation $U$, represented by the $2\times 2$ matrix $P$, is a reflection in the $x$-axis. \\

\begin{enumerate}[label=\textbf{(\alph*)}]
  \item Write down the matrix $P$. \marks{1} \\
\end{enumerate}

The transformation $V$, represented by the $2\times 2$ matrix $Q$, is a reflection in the line $y=x$. \\

\begin{enumerate}[label=\textbf{(\alph*)}, resume]
  \item Write down the matrix $Q$. \marks{1} \\
\end{enumerate}

Given that $U$ followed by $V$ is transformation $T$, which is represented by the matrix $R$ \\

\begin{enumerate}[label=\textbf{(\alph*)}, resume]
  \item express $R$ in terms of $P$ and $Q$, \marks{1}
  \item find the matrix $R$, \marks{2}
  \item give a full geometrical description of $T$ as a single transformation. \marks{2} \\
\end{enumerate}

\newpage

% ---- Question 6 ----
\questionitem

The matrix $\mathbf{M}$ is given by
\[
\mathbf{M} = \begin{pmatrix} 0 & 1 \\ 1 & 0 \end{pmatrix}\begin{pmatrix} 1 & 0 \\ 0 & k \end{pmatrix}
\]
where $k>0$ and $k \ne 1$. \\

\begin{enumerate}[label=\textbf{(\alph*)}]
  \item The matrix $\mathbf{M}$ represents a sequence of two geometrical transformations. State the type of each transformation, and make clear the order in which they are applied. \marks{2} \\
\end{enumerate}

The unit square in the $x$-$y$ plane is transformed by $\mathbf{M}$ onto rectangle $OABC$. \\

\begin{enumerate}[label=\textbf{(\alph*)}, resume]
  \item Find, in terms of $k$, the area of rectangle $OABC$ and the matrix which transforms $OABC$ onto the unit square. \marks{3} \\
  \item Show that the line through the origin with gradient $\frac{1}{\sqrt{k}}$ is invariant under the transformation in the $x$-$y$ plane represented by $\mathbf{M}$. \marks{3} \\
\end{enumerate}

\newpage

% ---- Question 7 ----
\questionitem

You are given the matrix $\mathbf{A} = \begin{pmatrix}0 & 1 & 0\\0 & 0 & 1\\1 & 0 & 0\end{pmatrix}$. \\

\begin{enumerate}[label=\textbf{(\alph*)}]
  \item Find $\mathbf{A}^5$. \marks{1} \\
  \item Describe the transformation that $\mathbf{A}$ represents. \marks{2} \\
\end{enumerate}

The matrix $\mathbf{B}$ represents a reflection in the plane $x = y$. \\

\begin{enumerate}[label=\textbf{(\alph*)}, resume]
  \item Write down the matrix $\mathbf{B}$. \marks{1} \\
\end{enumerate}

The point $Q$ has coordinates $(4, -1, 2)$. The point $Q'$ is the image of $Q$ under the transformation represented by $\mathbf{B}$. \\

\begin{enumerate}[label=\textbf{(\alph*)}, resume]
  \item Find the coordinates of $Q'$. \marks{1} \\
\end{enumerate}

\newpage

% ---- Question 8 ----
\questionitem

A right-angled triangle $T$ has vertices $A(-2,1)$, $B(1,1)$ and $C(1,5)$. When $T$ is transformed by the matrix
\[
P = \begin{pmatrix} 0 & -1 \\ 1 & 0 \end{pmatrix}
\]
the image is $T'$. \\

\begin{enumerate}[label=\textbf{(\alph*)}]
  \item Find the coordinates of the vertices of $T'$. \marks{2} \\
  \item Describe fully the transformation represented by $P$. \marks{2} \\
\end{enumerate}

The matrices
\[
Q = \begin{pmatrix} 2 & -1 \\ 0 & 3 \end{pmatrix}
\]
and
\[
R = \begin{pmatrix} 1 & 2 \\ 1 & 0 \end{pmatrix}
\]
represent two transformations. When $T$ is transformed by the matrix $QR$, the image is $T''$. \\

\begin{enumerate}[label=\textbf{(\alph*)}, resume]
  \item Find $QR$. \marks{2} \\
  \item Find the determinant of $QR$. \marks{2} \\
  \item Using your answer to part (d), find the area of $T''$. \marks{3} \\
\end{enumerate}

\newpage

% ---- Question 9 ----
\questionitem

Given that the matrix $\mathbf{M}$, where
\[
\mathbf{M} = \begin{pmatrix}
-\frac{\sqrt{2}}{2} & 0 & -\frac{\sqrt{2}}{2} \\
0 & 1 & 0 \\
\frac{\sqrt{2}}{2} & 0 & -\frac{\sqrt{2}}{2}
\end{pmatrix}
\]
represents a rotation through $\alpha^\circ$ about the $y$-axis in the positive sense given by the right-hand rule. \\

\begin{enumerate}[label=\textbf{(\alph*)}]
  \item determine the value of $\alpha$. \marks{1} \\
  \item Hence determine the smallest possible positive integer value of $k$ for which $\mathbf{M}^k = \mathbf{I}$. \marks{2} \\
\end{enumerate}

The $3\times 3$ matrix $\mathbf{N}$ represents a reflection in the plane with equation $x = 0$. \\

\begin{enumerate}[label=\textbf{(\alph*)}, resume]
  \item Write down the matrix $\mathbf{N}$. \marks{1} \\
\end{enumerate}

The point $A$ has coordinates $(1, 2, 3)$. \\

The point $B$ is the image of the point $A$ under the transformation represented by matrix $\mathbf{M}$ followed by the transformation represented by matrix $\mathbf{N}$. \\

\begin{enumerate}[label=\textbf{(\alph*)}, resume]
  \item Show that the coordinates of $B$ are $(2\sqrt{2}, 2, -\sqrt{2})$. \marks{2} \\
\end{enumerate}

Given that $O$ is the origin \\

\begin{enumerate}[label=\textbf{(\alph*)}, resume]
  \item show that, to $3$ significant figures, the size of angle $AOB$ is $79.4^\circ$. \marks{2} \\
  \item Hence determine the area of triangle $AOB$, giving your answer to $3$ significant figures. \marks{2} \\
\end{enumerate}

\newpage

% ---- Question 10 ----
\questionitem

Bilal has been asked to describe the transformation given by the matrix
\[
\begin{pmatrix}
\frac{\sqrt{3}}{2} & 0 & -\frac{1}{2} \\
0 & 1 & 0 \\
\frac{1}{2} & 0 & \frac{\sqrt{3}}{2}
\end{pmatrix}
\]

He writes his answer as follows: \\


\fbox{%
\parbox{\dimexpr\linewidth-2\fboxsep-2\fboxrule\relax}{%
The transformation is a rotation about the $y$-axis through an angle of $\theta$, where
\[
\cos \theta = \frac{\sqrt{3}}{2} \quad \text{and} \quad \sin \theta = \frac{1}{2}
\]
\[
\theta = 30^\circ
\]
}%
} \\[10pt]

Identify and correct the error in Bilal's work. \marks{2}

\newpage

% ---- Question 11 ----
\questionitem

Three non-singular $3 \times 3$ matrices, $\mathbf{A}$, $\mathbf{B}$ and $\mathbf{R}$, are such that
\[
\mathbf{R}\mathbf{A}=\mathbf{B}
\] \\
The matrix $\mathbf{R}$ represents a rotation about the $x$-axis through an angle $\theta$ and
\[
\mathbf{B}=\begin{bmatrix}
1 & 0 & 0 \\
0 & -\sin\theta & \cos\theta \\
0 & \cos\theta & \sin\theta
\end{bmatrix}
\] \\
\begin{enumerate}[label=\textbf{(\alph*)}]
  \item Show that $\mathbf{A}$ is independent of the value of $\theta$. \marks{3} \\
  \item Give a full description of the single transformation represented by the matrix $\mathbf{A}$. \marks{1} \\
\end{enumerate}
\end{enumerate}
\end{document}
